\documentclass[10pt,a4paper]{amsart}
\usepackage[margin=19mm]{geometry}
\usepackage{amsmath,amssymb,mathtools}
\usepackage[scr=rsfs,cal=euler]{mathalfa}
\usepackage{xfrac}
\usepackage[unicode,hidelinks,bookmarks=false]{hyperref}
\newcommand{\ie}{i.e.\ }
\newcommand{\cf}{cf.\ }
\newcommand{\resp}{respectively\ }
\newcommand{\Subsection}{\subsection}
\providecommand{\nobreakdash}{\nobreak\hbox{-}}
\setlength{\emergencystretch}{2em}
\multlinegap0pt
\usepackage{xcolor}
\usepackage{amssymb}
\usepackage{dsfont}
\usepackage{bm}
\usepackage{tikz}
\usepackage{tcolorbox}
\usepackage[shortlabels]{enumitem}
\usepackage{mathtools}
\newtheorem{lemma}{Lemma}
\newcommand{\CAL}[1]{\mathcal{#1}}
\title{A Covering Framework for Offline POMDPs Learning using Belief Space Metric}
\author{Youheng Zhu, Yiping Lu}
\date{Source: arXiv:2603.03191v1 (CC BY 4.0)}
\begin{document}
\maketitle

\noindent\emph{Source and licence.} A Covering Framework for Offline POMDPs Learning using Belief Space Metric. Version arXiv:2603.03191v1 (CC BY 4.0), distributed by arXiv under the Creative Commons Attribution 4.0 International licence, http://creativecommons.org/licenses/by/4.0/, as recorded in the arXiv licence metadata for this exact version and shown on the abstract page https://arxiv.org/abs/2603.03191v1. Full licence notice: \texttt{licenses/arxiv-2603.03191-CC-BY-4.0.txt}.\par\smallskip

\noindent\textit{Editorial scope.} Counted unit: the article's lemma lem:hoeffding\_result (source0.tex lines 1709--1714). The article's complete original proof of it is delivered with it (source0.tex lines 1715--1753, one \texttt{proof} environment of 39 lines). No mathematics is written, repaired or re-derived: the statement and the proof are the article's own lines, in the article's own notation, and no retained line is substituted or re-flowed.  No further source-local context is needed: every cross-reference of every retained line resolves inside the two delivered spans.  Nothing was omitted from any retained span, so the package carries no omission record.  No retained line contains a citation, so the wrapper carries no bibliography and no bibliography entry.  No retained line contains a figure environment, an image-inclusion command or drawing code, so no image is delivered and the package declares no asset dependency.  Editorial formatting only, in addition: an AMS \texttt{amsart} preamble that restates the article's own declarations and macros used by the retained lines, namely the environments \texttt{lemma} on separate counters, together with the standard packages those lines need.  The retained environments print in this document as Lemma 1 in order of appearance, whereas the article prints the same items with the numbering of its own full document.  The complete native source of the article as distributed in the arXiv e-print for this exact version is bundled unchanged as \texttt{sources/195660/source0.tex}.
\begin{lemma}\label{lem:hoeffding_result}
    With probability at least $1-\delta$, for $\forall f\in\mathcal{F}$,
    \begin{align}
        |\mathcal{E}_\phi(f,\pi)-\mathbb{E}_{d^{D}}[\mathcal{E}_\phi(f,\pi)]|\leq\sqrt{\frac{8R_{\rm max}^4}{n(1-\gamma)^4}\cdot\log\frac{2|\mathcal{F}|}{\delta}}
    \end{align}
\end{lemma}

\begin{proof}
For $\CAL{E}_\phi(f,\pi)$, we first estimate an upper bound on its absolute value. Since $f \in \CAL{F}$ is used to approximate a value function, its upper bound can be assumed to be no greater than $R{\rm max}/(1 - \gamma)$, i.e., the upper bound of the value function. Therefore, we can give a rough upper bound (possibly with a constant slack, which is acceptable since it’s only a constant):
\begin{align*}
0\leq \CAL{E}_\phi(f,\pi)\leq \frac{4R_{\rm max}^2}{(1-\gamma)^2}
\end{align*}

Thus, by Hoeffding’s inequality, for any $f \in \CAL{F}$, we have:
\begin{align}
&\ \Pr(|\mathcal{E}_\phi(f,\pi)-\mathbb{E}_{d^{D}}[\mathcal{E}_\phi(f,\pi)]|\geq t)\leq 2\exp\bigg(-\frac{2t^2n(1-\gamma)^4}{16R_{\rm max}^4}\bigg)\nonumber\\
        \to&\ \Pr(|\mathcal{E}_\phi(f,\pi)-\mathbb{E}_{d^{D}}[\mathcal{E}_\phi(f,\pi)]|> t)\leq 2\exp\bigg(-\frac{2t^2n(1-\gamma)^4}{16R_{\rm max}^4}\bigg)\label{ineq:hoeffding_target}
\end{align}

However, the goal of the proof is actually:
\begin{align}\label{equ:hoeffding_target_ultimate}
\Pr(\forall f\in\CAL{F}, |\mathcal{E}_\phi(f,\pi)-\mathbb{E}_{d^{D}}[\mathcal{E}_\phi(f,\pi)]|\leq t)
\end{align}

For such problems, a common approach is to use the union bound. Let the probability space be $(\Omega, \Sigma, \Pr)$, and define the events:
\begin{align*}
A:=&\ \{\omega\in\Omega:\forall f\in\CAL{F},|\mathcal{E}_\phi(f,\pi)-\mathbb{E}_{d^{D}}[\mathcal{E}_\phi(f,\pi)]|(\omega)\leq t\}\\
        B_f:=&\ \{\omega\in\Omega: |\mathcal{E}_\phi(f,\pi)-\mathbb{E}_{d^{D}}[\mathcal{E}_\phi(f,\pi)]|(\omega)> t\}
\end{align*}

Then:
\begin{align}
&\ \Pr(\forall f\in\CAL{F}, |\mathcal{E}_\phi(f,\pi)-\mathbb{E}_{d^{D}}[\mathcal{E}_\phi(f,\pi)]|\leq t)\nonumber\\
        =&\ \Pr(A)\nonumber
        = 1-\Pr(\Omega\backslash A)
        = 1-\Pr(\bigcup_{f\in\CAL{F}}B_f)\\
        \geq&\ 1-\sum_{f\in\CAL{F}}\Pr(B_f)
        \geq 1-2|\CAL{F}|\exp \bigg(-\frac{t^2n(1-\gamma)^4}{8R_{\rm max}^4}\bigg)
\end{align}

The second-to-last step uses the subadditivity of probability (countable subadditivity), and the final step applies inequality~\eqref{ineq:hoeffding_target}.

Let $\delta = 2|\CAL{F}|\exp \big(-{t^2n(1-\gamma)^4}/{8R_{\rm max}^4}\big)$, then solving for $t$ gives $t = \sqrt{\frac{8R_{\rm max}^4}{n(1-\gamma)^4}\cdot\log\frac{2|\CAL{F}|}{\delta}}$

Substituting this into~\eqref{equ:hoeffding_target_ultimate} completes the proof.
\end{proof}

\end{document}
